\documentclass[twoside,11pt]{article}

\usepackage{jmlr2e}
\usepackage{lastpage}

% Heading arguments: {volume}{year}{pages}{date submitted}{date published}{paper id}{author-full-names}
% The volume, page range, paper id and dates below are placeholders and will
% be filled in by the JMLR production editor upon acceptance.
\jmlrheading{27}{2026}{1-\pageref{LastPage}}{9/26}{}{26-0000}{Charles Poli, James Laidler, and Shreekanthadatta Seligar}

\ShortHeadings{Iguanas: Rule Generation and Selection on Polars}{Poli, Laidler, and Seligar}
\firstpageno{1}

\begin{document}

\title{Iguanas: A Python Library for Interpretable Decision Rule
Generation and Selection with ONNX Export}

\author{\name Charles Poli \email cpoli@paypal.com \\
       \addr PayPal, Inc.\\
       San Jose, CA, USA
       \AND
       \name James Laidler \email james.a.laidler@gmail.com \\
       \addr Independent Contributor
       \AND
       \name Shreekanthadatta Seligar \email seligar@paypal.com \\
       \addr PayPal, Inc.\\
       San Jose, CA, USA}

\editor{TBD}

\maketitle

\begin{abstract}
\texttt{Iguanas} is an open-source Python library for building rule-based
classifiers: models that predict by evaluating explicit, human-readable
conditions such as \texttt{(amount > 500) \& (account\_age\_days < 7)}. Built on
Polars, it provides a single API covering the full workflow --- generating
candidate rules from labelled data by extracting decision paths from
gradient-boosted tree ensembles, scoring and filtering candidates, composing a
final rule set under a choice of search strategies (including an exact
branch-and-bound search and a budgeted variant that constrains the fraction of
the population flagged), and exporting the result, including to ONNX, for
deployment. Beyond the core pipeline it includes modules for rule fairness
auditing, production monitoring, rule registries, and natural-language rule
explanation. On a standard imbalanced-classification benchmark under a
matched alert-rate budget, the budgeted composition strategy reaches recall
and precision comparable to existing rule learners while using an order of
magnitude fewer rule conditions. \texttt{Iguanas} is released under the
Apache-2.0 license, has been developed since 2021, and is used in production
at PayPal for transaction-risk rule development. Source code, documentation
and a reproducible benchmarking harness are publicly available.
\end{abstract}

\begin{keywords}
  interpretable machine learning, rule-based classification, imbalanced
  classification, Polars, ONNX
\end{keywords}

\section{Introduction and Statement of Need}

Interpretable models are preferable to post-hoc explanations of black-box
models wherever a prediction must be inspected, approved and defended rather
than merely accurate \citep{rudin2019,rudin2022}. Rule-based classifiers are
the most transparent members of that family: a rule can be read in a review
meeting, approved by a governance body, versioned, and audited when it
misfires.

Research software for rule learning is nonetheless fragmented. Individual
induction algorithms are well served --- \texttt{RuleFit} \citep{friedman2008},
\texttt{SkopeRules} \citep{skope}, certifiably optimal rule lists
\citep{angelino2017}, sparse decision trees \citep{lin2020}, and the
\texttt{imodels} package \citep{singh2021}, which collects several of these
including FIGS \citep{tan2025}. Each of these implements a particular
\emph{induction} algorithm. Researchers who need the surrounding workflow ---
produce a large and deliberately diverse candidate pool, score and de-duplicate
it, select a subset under an explicit operating constraint, then audit and
deploy the result --- typically assemble it themselves, which makes studies
difficult to reproduce and comparisons difficult to control. \texttt{Iguanas}
provides that workflow as a single library, with each stage separately
configurable so that generation, filtering and selection can be varied
independently in an experiment.

Two capabilities are less common in comparable research tooling. First, rule
sets can be exported to ONNX \citep{onnx}, so a rule set developed in Python
can be served by any ONNX runtime --- useful for research on models intended
for environments where a Python interpreter is unavailable, or where the
scoring artefact must be frozen and signed. Second, the \texttt{rule\_formatting}
module translates rules stated over \emph{preprocessed} features back into the
original columns, so that rules learned after imputation, encoding or binning
remain human-readable rather than expressed in terms of transformed variables.

\section{Software Description}

\texttt{Iguanas} covers four pipeline stages in one API, plus a set of
auxiliary modules for the surrounding operational lifecycle.

\textbf{Generation.} Rules are generated by extracting decision paths from
fitted XGBoost/LightGBM tree ensembles \citep{chen2016,ke2017}, guided by
monotone constraints, with a grid over class weights and sample-weight
transformations used to diversify the candidate pool.

\textbf{Evaluation and filtering.} Candidates are scored with standard
classification metrics and with coverage-aware rule-quality measures (lift,
weighted relative accuracy, Laplace and m-estimate), then filtered on
performance, feature overlap and pairwise correlation.

\textbf{Composition.} A final rule set is composed from the surviving
candidates by a choice of search strategies: greedy, cumulative, beam search,
an exact branch-and-bound search, exhaustive enumeration, and a budgeted
variant that greedily maximises marginal coverage under an alert-rate
constraint (a $(1-1/e)$-approximate solution to budgeted maximum coverage).
The branch-and-bound search prunes using an admissible upper bound on recall,
precision, accuracy or $F_\beta$ at each node, obtained from the monotonicity
of true and false positives under the rule-combination operator: it returns
the same optimum as exhaustive enumeration while expanding a fraction of the
nodes.

\textbf{Deployment.} Each rule is compiled once into a Polars \citep{polars}
expression and cached, so evaluation is columnar and multi-threaded across the
full rule set. Rule sets can additionally be exported to ONNX for deployment
outside Python. Grid search over generation parameters is parallelised with
joblib threading; execution is single-node (no multiprocessing, distributed
backend or GPU path).

\textbf{Auxiliary modules.} Beyond the core pipeline, \texttt{Iguanas}
provides post-hoc rule fairness auditing (subgroup metrics and disparate
impact ratio), production monitoring (metric-delta comparison across periods),
a rule registry for versioned storage, and natural-language rule explanation.
A scikit-learn-compatible \texttt{RulesetClassifier} estimator wraps
generation, filtering and combination into a single \texttt{fit}/\texttt{predict}
object for users who want the default pipeline rather than its individual
stages.

\section{Comparison to Related Software}

Unlike single-algorithm libraries (\texttt{RuleFit}, \texttt{SkopeRules},
\texttt{imodels}), \texttt{Iguanas} treats generation, evaluation, filtering
and composition as independently configurable stages of one pipeline, rather
than a fixed algorithm. This makes it possible to hold the induction method
fixed while varying selection strategy (or vice versa) in a controlled
experiment. The library also differs from these tools in providing an
ONNX export path and reversible preprocessing-aware rule formatting, and in
providing a budgeted composition strategy that targets a specific
operating point (fraction of the population flagged) directly, rather than
requiring the user to search a score threshold.

The repository includes a benchmarking harness that runs \texttt{Iguanas}
alongside several external rule learners (RIPPER, Bayesian Rule Lists, FIGS,
decision trees, RuleFit, SkopeRules) over a registry of public imbalanced
classification datasets, under a nested generate/select/test protocol with a
runtime leakage guard, so that generation, selection and evaluation are done
on disjoint splits. It records candidate-pool size, fitted-tree counts,
timing and rule-set complexity (number of conditions) alongside quality
metrics at a matched alert-rate budget --- the fraction of the population the
rule set is allowed to flag, which is the operating constraint that matters
in practice and the one score-based learners can otherwise satisfy trivially
by sliding a threshold.

Table~\ref{tab:mammography} illustrates the comparison on \texttt{mammography}
\citep{mammography}, a standard imbalanced-classification benchmark, at two
alert-rate budgets. \texttt{Iguanas} does not dominate on raw recall, but it
reaches recall within 0.125 of \texttt{RuleFit} at the 5\% budget, and
exceeds \texttt{SkopeRules} at matched precision, using an order of magnitude
fewer rule conditions than either. At the 1\% budget it is essentially tied
with \texttt{SkopeRules} on recall and precision at an 8-fold reduction in
complexity. This is the defensible claim: comparable predictive performance
at much lower rule-set complexity, not raw dominance. This is one dataset,
shown as a worked illustration; the harness is designed to run the same
protocol across the full dataset registry, and broader results are
maintained alongside the repository as the comparison is extended.

\begin{table}[t]
\centering
\caption{Matched alert-rate-budget comparison on \texttt{mammography}, nested
generate/select/test protocol. Alert rate (ar) is the fraction of the test
set flagged; conditions is the total number of atomic conditions in the
final rule set (lower is simpler).}
\label{tab:mammography}
\begin{tabular}{lccccc}
\hline
\textbf{Budget} & \textbf{Model} & \textbf{Recall} & \textbf{Precision} & \textbf{ar} & \textbf{Conditions} \\
\hline
5\%  & \texttt{Iguanas}    & 0.665 & 0.476 & 0.032 & 9 \\
5\%  & \texttt{SkopeRules} & 0.550 & 0.473 & 0.031 & 39 \\
5\%  & \texttt{RuleFit}    & 0.790 & 0.354 & --    & 62 \\
5\%  & GBM ceiling         & 0.785 & 0.466 & --    & 13139 \\
\hline
1\%  & \texttt{Iguanas}    & 0.360 & 0.819 & --    & 5 \\
1\%  & \texttt{SkopeRules} & 0.375 & 0.860 & --    & 39 \\
\hline
\end{tabular}
\end{table}

\section{Availability, Documentation, and Community}

\texttt{Iguanas} has been developed since 2021 and is used in production at
PayPal for transaction-risk rule development. It is released under the
Apache-2.0 license. The source code is hosted at
\url{https://github.com/paypal/iguanas}, with issue tracking and pull-request
review conducted there. Releases are published to PyPI
(\texttt{pip install iguanas}), with optional extras for ONNX export and
LightGBM support. Documentation, including a full API reference and worked
notebooks, is published at \url{https://paypal.github.io/Iguanas/}. The test
suite (over 700 tests) runs under continuous integration across five Python
versions on every push to the main branch.

\acks{We thank the PSP Data Team at PayPal, and Jem Day and Aditya Kapadiya
for contributions to the codebase. This work was supported by PayPal, Inc.
The authors declare no competing interests beyond their employment.}

\vskip 0.2in
\bibliography{paper}

\end{document}
